\documentclass[a4paper,11pt]{article}
\usepackage[T1]{fontenc}
\usepackage[utf8]{inputenc}
\usepackage[ngerman]{babel}
\usepackage{amsmath,amssymb}

\title{EABC-Volumenkonstanz und dyadische Gitterteilung}
\date{}

\begin{document}
\maketitle

Wir betrachten das lokale Einheitsvolumen \(V_{\mathcal M}\) eines
arithmetisch-geometrischen Quelltraegers. Dieses Volumen werde durch vier
orthogonale EABC-Koordinatenrichtungen
\[
\hat{\mathbf e},\hat{\mathbf a},\hat{\mathbf b},\hat{\mathbf c}
\]
aufgespannt. Die vier Richtungen entsprechen den vier invertierbaren
Restklassen modulo \(12\),
\[
E \equiv 1,\qquad A \equiv 5,\qquad B \equiv 7,\qquad C \equiv 11 \pmod{12}.
\]

Die dyadische Zerlegung des lokalen Einheitsvolumens erfolgt entlang eines
periodischen Vierphasenpfades. Wir definieren die vektorielle
Amplitudenreihe
\[
\boldsymbol{\Xi}
=
\sum_{k=0}^{\infty}
\left(
\frac{1}{2^{4k+1}}\hat{\mathbf e}
+
\frac{1}{2^{4k+2}}\hat{\mathbf a}
+
\frac{1}{2^{4k+3}}\hat{\mathbf b}
+
\frac{1}{2^{4k+4}}\hat{\mathbf c}
\right).
\]

Da nach vier Schritten derselbe EABC-Zyklus erneut beginnt, besitzt jede
Koordinatenrichtung denselben geometrischen Skalenfaktor
\[
q=\frac{1}{2^4}=\frac{1}{16}.
\]

Damit ergeben sich die vier Komponenten
\[
\Xi_e
=
\sum_{k=0}^{\infty}\frac{1}{2^{4k+1}}
=
\frac{1}{2}\cdot\frac{1}{1-\frac{1}{16}}
=
\frac{8}{15},
\]
\[
\Xi_a
=
\sum_{k=0}^{\infty}\frac{1}{2^{4k+2}}
=
\frac{1}{4}\cdot\frac{1}{1-\frac{1}{16}}
=
\frac{4}{15},
\]
\[
\Xi_b
=
\sum_{k=0}^{\infty}\frac{1}{2^{4k+3}}
=
\frac{1}{8}\cdot\frac{1}{1-\frac{1}{16}}
=
\frac{2}{15},
\]
und
\[
\Xi_c
=
\sum_{k=0}^{\infty}\frac{1}{2^{4k+4}}
=
\frac{1}{16}\cdot\frac{1}{1-\frac{1}{16}}
=
\frac{1}{15}.
\]

Somit besitzt die EABC-Zerlegung die normierte Gesamtmasse
\[
\sum_{\mu\in\{e,a,b,c\}}\Xi_\mu
=
\frac{8}{15}
+
\frac{4}{15}
+
\frac{2}{15}
+
\frac{1}{15}
=
1.
\]

Die Reihe definiert daher eine vollstaendige additive Zerlegung des lokalen
Einheitsvolumens. Es entstehen weder Luecken noch Ueberdeckungen im Sinne
der \(L^1\)-Norm.

Die zugehoerige gewichtete Richtungsabbildung lautet
\[
\boldsymbol{\Xi}
=
\frac{8}{15}\hat{\mathbf e}
+
\frac{4}{15}\hat{\mathbf a}
+
\frac{2}{15}\hat{\mathbf b}
+
\frac{1}{15}\hat{\mathbf c}.
\]

Da die vier Basisrichtungen orthogonal sind, besitzt dieser Vektor die
euklidische Norm
\[
\lVert\boldsymbol{\Xi}\rVert_2^2
=
\frac{64+16+4+1}{225}
=
\frac{17}{45}.
\]

Die normierte SLERP-Richtung ist daher nicht \(\boldsymbol{\Xi}\) selbst,
sondern
\[
\widehat{\boldsymbol{\Xi}}
=
\frac{\boldsymbol{\Xi}}{\lVert\boldsymbol{\Xi}\rVert_2}
=
\sqrt{\frac{45}{17}}\,\boldsymbol{\Xi}.
\]

Damit trennt das Modell sauber zwischen additiver Volumenkonstanz und
geodaetischer Richtungsnormierung: Die dyadische EABC-Zerlegung erhaelt die
Gesamtmasse exakt, waehrend die SLERP-Normierung die zugehoerige Richtung auf
der Einheitssphaere des quaternionischen Koordinatenraums fixiert.

Die Zahl \(15\) tritt in diesem Zusammenhang nicht zufaellig auf. Sie ist die
Gesamtmasse der vier dyadischen Gewichte
\[
8+4+2+1=15
\]
und zugleich die Anzahl der nichtleeren Teilmengen eines Vierer-Systems. Im
EABC-Modell kann sie daher als kombinatorischer Nenner der vollstaendigen
Vierphasen-Zerlegung verstanden werden. Die vier Komponenten \((8,4,2,1)\)
bilden dabei eine hierarchische Gewichtung der EABC-Achsen, ohne die
Gesamtmasse des lokalen Quelltraegers zu veraendern.

\paragraph{Lemma.}
Die dyadische Vierphasen-Zerlegung des EABC-Koordinatenraums liefert eine
eindeutig normierte additive Zerlegung des lokalen Einheitsvolumens:
\[
V_{\mathcal M}
=
\Xi_e+\Xi_a+\Xi_b+\Xi_c
=
1.
\]

Die zugehoerigen Koordinatengewichte sind
\[
(\Xi_e,\Xi_a,\Xi_b,\Xi_c)
=
\left(
\frac{8}{15},
\frac{4}{15},
\frac{2}{15},
\frac{1}{15}
\right).
\]

Damit ist die unendliche dyadische Zerlegung im EABC-Raum konvergent,
vollstaendig und volumenstabil im additiven Sinn. Eine zusaetzliche
geodaetische Normierung erzeugt daraus eine Richtung auf der quaternionischen
Einheitssphaere.

\end{document}
